\documentclass[12pt]{article}
\usepackage[utf8]{inputenc}
\usepackage[T1]{fontenc}
\usepackage[margin=1in]{geometry}
\usepackage{setspace}
\usepackage{xcolor}
\usepackage[colorlinks=true, urlcolor=blue!60!black, linkcolor=blue!60!black]{hyperref}
\setlength{\parindent}{0em}
\setlength{\parskip}{0.6em}
\renewcommand{\baselinestretch}{1.0}
\begin{document}
\begin{center}
{\LARGE \textit{Missouri v. Holland}, 252 U.S. 416 (1920)\par}
\vspace{0.5em}
{\large Case Brief --- Two-Page Edition\par}
\end{center}

\section{Memory Jogger}
Congress may enact legislation necessary and proper to execute a valid treaty, even on a subject it could not otherwise regulate, without violating the Tenth Amendment (\href{https://supreme.justia.com/cases/federal/us/252/416/}{Justia}).

\section{Detailed Case Facts}
After courts struck down a 1913 federal bird-hunting law, the United States and Great Britain (for Canada) signed the 1916 Migratory Bird Treaty to save endangered migratory species, and Congress implemented it in the 1918 Migratory Bird Treaty Act.

\section{Procedural History}
Missouri sued to enjoin federal game warden Ray Holland from enforcing the Act. The district court dismissed, upholding the Act, and Missouri appealed directly.

\section{Judicial Votes}
Affirmed, 7--2. Holmes for the Court (White, McKenna, Day, McReynolds, Brandeis, Clarke); Van Devanter and Pitney dissented without opinion.

\section{Holding}
The treaty is a valid exercise of the delegated treaty power, and the Act is a necessary and proper means of executing it. A valid treaty and its implementing law are not limited to what Congress could do alone under Article~I.

\section{Analysis of Opinions}
\textbf{Holmes:} The Tenth Amendment reserves only undelegated powers, and treaties are supreme when made under the authority of the United States. The Constitution created a nation whose development its framers could not foresee. No ``prohibitory words'' barred the treaty, only ``some invisible radiation'' from the Tenth Amendment. Missouri's title claim failed, and protecting migrating birds required national action. \textbf{Later:} In \textit{Bond} (2014), Scalia, Thomas, and Alito questioned \textit{Holland}'s reach.

\section{Examples: Future Applications}
\textbf{Same-Side:} A treaty-based ban on destroying an endangered migrating butterfly's habitat is valid. \textbf{Opposite-Side:} A treaty-based ban on insulting religion violates the First Amendment (\textit{Reid v. Covert}). \textbf{Fence-Sitter:} A treaty-based federal ban on school corporal punishment tests the \textit{Bond} concurrences' limits.

\section{Critique}
Rosenkranz argues Holmes misread Article~VI and the Necessary and Proper Clause. Others note the collective-action rationale has no stopping point against pretextual treaties.

\section{Key Quotations}
\begin{quote}
``It was enough for them to realize or to hope that they had created an organism; it has taken a century and has cost their successors much sweat and blood to prove that they created a nation. The case before us must be considered in the light of our whole experience, and not merely in that of what was said a hundred years ago.''
\end{quote}
\noindent --- Justice Holmes, 252 U.S. at 433.

\end{document}
